\section{\texorpdfstring{Injective and projective resolutions:\allowbreak{} signs,\allowbreak{} relative lifts,\allowbreak{} and fixed outer terms}{Injective and projective   relative  and fixed outer terms}}\label{proofs-6209-7089-injective-and-projective-resolutions-signs-relative-lifts-and-fixed-outer-terms}

This editorial review covers original Derived Categories 6209–\allowbreak{}7089 at authority commit a04446e5\allowbreak{}7ec1fbc2\allowbreak{}52a871af\allowbreak{}cec7752f\allowbreak{}b2807b14. Earlier corrections are checked against that source and the frozen current chapter. The two strengthenings below remain separate from the source body. They are consequences of the reviewed constructions; no novelty is asserted. Cohomological shifts have differential d\_\allowbreak{}\{X[1]\allowbreak{}\}\textasciicircum{}n\allowbreak{}=-d\_\allowbreak{}X\textasciicircum{}\{n\allowbreak{}+1\},\allowbreak{} and a homotopy from a map f to zero satisfies f\textasciicircum{}n\allowbreak{}=d\_\allowbreak{}Y\textasciicircum{}\{n-1\}h\textasciicircum{}n\allowbreak{}+h\textasciicircum{}\{n\allowbreak{}+1\}d\_\allowbreak{}X\textasciicircum{}n.

\subsection{1. Existence and the original bounds}\label{proofs-6209-7089-existence-and-the-original-bounds}

An injective resolution of an object A is the source\textquotesingle s complex I with I\textasciicircum{}n\allowbreak{}=0 for n<0,\allowbreak{} injective terms,\allowbreak{} A identified with ker(d\_\allowbreak{}I\textasciicircum{}0)\allowbreak{},\allowbreak{} and H\textasciicircum{}n(I)\allowbreak{}\allowbreak{}=0 for n\textgreater0. Thus A[0]\allowbreak{} \allowbreak{}-> I is a quasi-isomorphism. A resolution of a general complex K requires I to be bounded below,\allowbreak{} so its existence forces H\textasciicircum{}n(K)\allowbreak{}\allowbreak{}=0 in all sufficiently small degrees.

Conversely,\allowbreak{} if H\textasciicircum{}n(K)\allowbreak{}\allowbreak{}=0 for n<a,\allowbreak{} use the actual canonical truncation q:\allowbreak{}K \allowbreak{}-> tau\_\allowbreak{}\{>\allowbreak{}=a\}K. Its degree-a term is coker(d\_\allowbreak{}K\textasciicircum{}\{a-1\})\allowbreak{},\allowbreak{} its terms above a are K\textasciicircum{}n,\allowbreak{} and its terms below a are zero. The map q induces the identity identifications on H\^{}n for n>\allowbreak{}=a and is an isomorphism on the zero cohomology below a. It is therefore a quasi-isomorphism.

If the category has enough injectives,\allowbreak{} embed A into I\textasciicircum{}0,\allowbreak{} then embed I\textasciicircum{}0/\allowbreak{}A into I\textasciicircum{}1,\allowbreak{} then coker(d\_\allowbreak{}I\textasciicircum{}0)\allowbreak{} into I\textasciicircum{}2,\allowbreak{} and continue. The differential is the quotient map followed by the chosen monomorphism. Its kernel is the preceding image,\allowbreak{} including ker(d\_\allowbreak{}I\textasciicircum{}0)\allowbreak{}\allowbreak{}=A; successive differentials compose to zero. This proves the object resolution with all its original terms. DERIVED-RECON-036 reconciles P02-E0041 and French DERIVED-033 with MC-STK-ERR-0844:\allowbreak{} the quotient at original 6297 is I\textasciicircum{}0/\allowbreak{}A,\allowbreak{} not the undefined I\_\allowbreak{}0/\allowbreak{}A.

For a complex whose terms are zero below a,\allowbreak{} apply the source\textquotesingle s subcategory-right-resolution lemma,\allowbreak{} original 5386–\allowbreak{}5405,\allowbreak{} to the class of injective objects. It produces a termwise monomorphism K \allowbreak{}-> I,\allowbreak{} an isomorphism on cohomology,\allowbreak{} and I\textasciicircum{}n\allowbreak{}=0 for n\textless a. The construction and its exact pushout maps are proved in section 3 of PROOFS\_\allowbreak{}5091\_\allowbreak{}5900.md,\allowbreak{} by dualizing its fully written pullback induction. This application satisfies both hypotheses:\allowbreak{} zero is injective,\allowbreak{} and every object embeds into an injective. For K with only the cohomological bound,\allowbreak{} compose the quasi-isomorphism q above with the resolution of tau\_\allowbreak{}\{>\allowbreak{}=a\}K. That composite need not be a termwise monomorphism; no such assertion is made in the cohomologically bounded case.

\subsection{2. The null-homotopy induction and its corrected indices}\label{proofs-6209-7089-the-null-homotopy-induction-and-its-corrected-indices}

Let K be acyclic,\allowbreak{} let every I\^{}n be injective,\allowbreak{} and suppose I\textasciicircum{}n\allowbreak{}=0 for n\textless a. For a chain map alpha:\allowbreak{}K \allowbreak{}-> I set h\textasciicircum{}j\allowbreak{}=0 for j<\allowbreak{}=a. Suppose h has been defined through h\^{}n with

\[
 \alpha^j=d_I^{j-1}h^j+h^{j+1}d_K^j\quad(j<n).
\]

The next defect is b\textasciicircum{}n\allowbreak{}=alpha\textasciicircum{}n-d\_\allowbreak{}I\textasciicircum{}\{n-1\}h\textasciicircum{}n:\allowbreak{}K\textasciicircum{}n \allowbreak{}-> I\^{}n. With every domain explicit,\allowbreak{} the source calculation is

\[
\begin{aligned}
 b^n d_K^{n-1}
 &=\alpha^n d_K^{n-1}-d_I^{n-1}h^n d_K^{n-1}\\
 &=d_I^{n-1}\alpha^{n-1}-d_I^{n-1}h^n d_K^{n-1}\\
 &=d_I^{n-1}(d_I^{n-2}h^{n-1}+h^n d_K^{n-1})
                      -d_I^{n-1}h^n d_K^{n-1}=0.          
\end{aligned}
\tag{2.1}\]

Thus b\^{}n kills Im(d\_\allowbreak{}K\textasciicircum{}\{n-1\})\allowbreak{}\allowbreak{}=ker(d\_\allowbreak{}K\textasciicircum{}n)\allowbreak{}. The universal property of the cokernel identifies its induced map with a map Im(d\_\allowbreak{}K\textasciicircum{}n)\allowbreak{} \allowbreak{}-> I\^{}n. Extend it along Im(d\_\allowbreak{}K\textasciicircum{}n)\allowbreak{} \allowbreak{}-> K\textasciicircum{}\{n\allowbreak{}+1\},\allowbreak{} using injectivity of I\textasciicircum{}n,\allowbreak{} to h\textasciicircum{}\{n\allowbreak{}+1\}:\allowbreak{}K\textasciicircum{}\{n\allowbreak{}+1\} \allowbreak{}-> I\^{}n. The resulting equation is precisely alpha\textasciicircum{}n\allowbreak{}=d\_\allowbreak{}I\textasciicircum{}\{n-1\}h\textasciicircum{}n\allowbreak{}+h\textasciicircum{}\{n\allowbreak{}+1\}d\_\allowbreak{}K\textasciicircum{}n. Induction proves the null-homotopy in all degrees without an assumption on the lower terms of K.

DERIVED-RECON-037 reconciles P02-E0042 and French DERIVED-034 with all six operations in MC-STK-ERR-0845. Five occurrences of h\^{}\{n-1\} outside the previous-degree identity must be h\textasciicircum{}n,\allowbreak{} and the first alpha in original 6328 must be alpha\^{}n. The occurrence h\^{}\{n-1\} inside d\_\allowbreak{}I\textasciicircum{}\{n-2\}h\textasciicircum{}\{n-1\} in (2.1)\allowbreak{} remains unchanged. The current source implements this exact selective correction.

For a quasi-isomorphism alpha:\allowbreak{}K \allowbreak{}-> L,\allowbreak{} its original cone C(alpha)\allowbreak{} is acyclic. The distinguished homotopy-category triangle is (K,\allowbreak{}L,\allowbreak{}C(alpha)\allowbreak{},\allowbreak{}alpha,\allowbreak{}i,\allowbreak{}-p)\allowbreak{}. Applying Hom\_\allowbreak{}K(-,\allowbreak{}I)\allowbreak{} gives the exact segment

\[
 \operatorname{Hom}_K(C(\alpha),I)\longrightarrow
 \operatorname{Hom}_K(L,I)\xrightarrow{\alpha^*}
 \operatorname{Hom}_K(K,I)\longrightarrow
 \operatorname{Hom}_K(C(\alpha)[-1],I).
\]

Both outer groups vanish by the just proved induction,\allowbreak{} including the shifted acyclic source. Hence alpha\^{}* is an isomorphism. This proves existence and uniqueness of a lift up to homotopy independently of the direct proofs that follow.

\subsection{3. Strict lifting and the remaining sign error}\label{proofs-6209-7089-strict-lifting-and-the-remaining-sign-error}

Suppose alpha:\allowbreak{}K \allowbreak{}-> L is a termwise monomorphic quasi-isomorphism and gamma:\allowbreak{}K \allowbreak{}-> I is given. We construct beta:\allowbreak{}L \allowbreak{}-> I with beta alpha\allowbreak{}=gamma. Set beta\textasciicircum{}j\allowbreak{}=0 in degrees j\textless a. Assume the construction through degree n-1 is compatible with the differentials. To prescribe beta\^{}n on Im(d\_\allowbreak{}L\textasciicircum{}\{n-1\})\allowbreak{},\allowbreak{} first show that d\_\allowbreak{}I\textasciicircum{}\{n-1\}beta\textasciicircum{}\{n-1\} kills ker(d\_\allowbreak{}L\textasciicircum{}\{n-1\})\allowbreak{}.

It kills Im(d\_\allowbreak{}L\textasciicircum{}\{n-2\})\allowbreak{} by the already established chain equation. Its restriction to ker(d\_\allowbreak{}L\textasciicircum{}\{n-1\})\allowbreak{} therefore induces a map H\textasciicircum{}\{n-1\}(L)\allowbreak{} \allowbreak{}-> I\^{}n. Composing with H\textasciicircum{}\{n-1\}(alpha)\allowbreak{} gives zero:\allowbreak{} on K-cycles it is d\_\allowbreak{}I\textasciicircum{}\{n-1\}gamma\textasciicircum{}\{n-1\}\allowbreak{}=gamma\textasciicircum{}n d\_\allowbreak{}K\textasciicircum{}\{n-1\}\allowbreak{}=0. Since H\textasciicircum{}\{n-1\}(alpha)\allowbreak{} is an isomorphism,\allowbreak{} the induced map is zero. This proves the required factorization through Im(d\_\allowbreak{}L\textasciicircum{}\{n-1\})\allowbreak{}. At the starting degree it also follows immediately from I\textasciicircum{}\{a-1\}\allowbreak{}=0.

The prescription on alpha(K\textasciicircum{}n)\allowbreak{} is gamma\^{}n. The intersection identity needed to glue it to the prescription on Im(d\_\allowbreak{}L\textasciicircum{}\{n-1\})\allowbreak{} is

\[
 \alpha(\operatorname{Im}d_K^{n-1})
       =\alpha(K^n)\cap\operatorname{Im}d_L^{n-1}.         \tag{3.1}
\]

The forward inclusion follows from the chain equation. For the reverse,\allowbreak{} pull back the intersection along the monomorphism alpha\^{}n. Its objects lie in ker(d\_\allowbreak{}K\textasciicircum{}n)\allowbreak{},\allowbreak{} because alpha\textasciicircum{}\{n\allowbreak{}+1\} is monic; their images in H\textasciicircum{}n(L)\allowbreak{} vanish. Injectivity of H\textasciicircum{}n(alpha)\allowbreak{} forces their images in H\textasciicircum{}n(K)\allowbreak{} to vanish,\allowbreak{} so this pullback is contained in Im(d\_\allowbreak{}K\textasciicircum{}\{n-1\})\allowbreak{}. This proves the subobject identity (3.1)\allowbreak{} in the abelian category. On that intersection the two maps agree by beta\^{}\{n-1\}alpha\^{}\{n-1\} \allowbreak{}=gamma\textasciicircum{}\{n-1\} and the chain equation for gamma. They consequently give a map from the sum of the two subobjects to I\^{}n. Injectivity extends it to all of L\^{}n. Induction proves strict lifting. This also supplies the implicit well-definedness argument in the source\textquotesingle s direct proof; no additional source-body reconstruction is required.

For general alpha use exactly the source\textquotesingle s factorization at 2622–\allowbreak{}2724:\allowbreak{}

\[
 \widetilde L^n=L^n\oplus K^n\oplus K^{n+1},\qquad
 d_{\widetilde L}^n=
 \begin{pmatrix}d_L^n&0&0\\0&d_K^n&1\\0&0&-d_K^{n+1}\end{pmatrix},
\] \[
 \widetilde\alpha=(\alpha,1,0)^t,\qquad
 \pi=(1,0,0),\qquad s=(1,0,0)^t.                         \tag{3.2}
\]

Here pi tilde-alpha\allowbreak{}=alpha,\allowbreak{} pi s\allowbreak{}=1,\allowbreak{} and the original homotopy H(l,\allowbreak{}k,\allowbreak{}k\textasciicircum{}\allowbreak{}+)\allowbreak{}\allowbreak{}=(0,\allowbreak{}0,\allowbreak{}k)\allowbreak{} satisfies

\[
 (dH+Hd)(l,k,k^+)=(0,k,-dk)+(0,0,dk+k^+)=(0,k,k^+)
                        =(1-s\pi)(l,k,k^+).             \tag{3.3}
\]

Thus tilde-alpha is a split monomorphic quasi-isomorphism. A strict lift tilde-beta gives beta\allowbreak{}=tilde-beta s,\allowbreak{} with beta alpha homotopic to gamma by (3.3)\allowbreak{}. DERIVED-RECON-038 reconciles both occurrences,\allowbreak{} 6412 and 6502,\allowbreak{} in P02-E0043 and French DERIVED-035 with the two MC-STK-ERR-0846 operations restoring the complex symbol K\^{}bullet.

For the source\textquotesingle s direct uniqueness argument,\allowbreak{} after this replacement write L\textasciicircum{}n\allowbreak{}=K\textasciicircum{}n direct-sum M\textasciicircum{}n,\allowbreak{} where M is acyclic. The given homotopies satisfy

\[
 \gamma^n-\beta_i^n|_{K^n}
       =d_I^{n-1}h_i^n+h_i^{n+1}d_K^n.
\]

Subtracting the two displayed equations in the actual order gives

\[
 \beta_1^n|_{K^n}-\beta_2^n|_{K^n}
 =d_I^{n-1}(h_2^n-h_1^n)
                  +(h_2^{n+1}-h_1^{n+1})d_K^n.          \tag{3.4}
\]

\textbf{DERIVED-NEW-024:\allowbreak{}} original 6524 and 6527,\allowbreak{} still unchanged in the current chapter,\allowbreak{} use h\_\allowbreak{}1-h\_\allowbreak{}2 instead. Correct both to h\_\allowbreak{}2-h\_\allowbreak{}1. The original maps beta\_\allowbreak{}1,\allowbreak{}beta\_\allowbreak{}2,\allowbreak{}gamma and the original convention for each h\_\allowbreak{}i are retained; reversing a convention elsewhere would conceal the actual error.

Extend h\_\allowbreak{}2\textasciicircum{}n-h\_\allowbreak{}1\textasciicircum{}n by zero on the chosen M\^{}n summand and call it h\^{}n. Then

\[
 r^n=\beta_1^n-\beta_2^n
                -(d_I^{n-1}h^n+h^{n+1}d_L^n)            \tag{3.5}
\]

is a chain map and is zero on K. It factors through the quotient q:\allowbreak{}L \allowbreak{}-> M. Write r\allowbreak{}=tq. Section 2 gives t\allowbreak{}=d\_\allowbreak{}I k\allowbreak{}+k d\_\allowbreak{}M,\allowbreak{} so beta\_\allowbreak{}1-beta\_\allowbreak{}2\allowbreak{}=d\_\allowbreak{}I(h\allowbreak{}+kq)\allowbreak{}\allowbreak{}+(h\allowbreak{}+kq)\allowbreak{}d\_\allowbreak{}L. This is the required homotopy. DERIVED-RECON-039 confirms the separate MC-STK-ERR-0847 parenthesis repair reported by P02-E0044 and French DERIVED-036. It is retained in (3.5)\allowbreak{}; it did not correct the independent sign in (3.4)\allowbreak{}.

For a counterexample to the uncorrected intermediate formula,\allowbreak{} work in abelian groups and take K\allowbreak{}=L\allowbreak{}=I\allowbreak{}=(Q --1-\allowbreak{}-> Q)\allowbreak{} in degrees 0,\allowbreak{}1. All terms of I are injective. Put alpha\allowbreak{}=1,\allowbreak{} gamma\allowbreak{}=0,\allowbreak{} beta\_\allowbreak{}1\allowbreak{}=1,\allowbreak{} beta\_\allowbreak{}2\allowbreak{}=0,\allowbreak{} h\_\allowbreak{}1\textasciicircum{}1\allowbreak{}=-1 and h\_\allowbreak{}2\textasciicircum{}1\allowbreak{}=0,\allowbreak{} with all other homotopy components zero. These satisfy both given homotopy equations. The source\textquotesingle s h\_\allowbreak{}1-h\_\allowbreak{}2 has coboundary -1,\allowbreak{} whereas beta\_\allowbreak{}1-beta\_\allowbreak{}2\allowbreak{}=1. Its proposed residual is 2 times the identity and is not zero on K\allowbreak{}=L. The corrected h\_\allowbreak{}2-h\_\allowbreak{}1 has coboundary 1 and gives zero residual. Thus the formula needs correction,\allowbreak{} while the uniqueness theorem stays valid.

Finally,\allowbreak{} every derived arrow L \allowbreak{}-> I is a roof gamma alpha\^{}\{-1\}. Lifting gamma produces a chain map beta representing it. If the class of a chain map beta becomes zero,\allowbreak{} some quasi-isomorphism alpha makes beta alpha null-homotopic; the injectivity of alpha\^{}* proved in section 2 makes beta null-homotopic. Consequently

\[
 \operatorname{Hom}_{K(A)}(L,I)
       \xrightarrow{\sim}\operatorname{Hom}_{D(A)}(L,I).  \tag{3.6}
\]

\textbf{DERIVED-NEW-025:\allowbreak{}} the introductions of this lemma at original 6543 and its projective dual at 6824 each omit the article before ``bounded''. Insert ``a'' in those two sentences. This is a copyedit,\allowbreak{} with no changed bounds or mathematical hypotheses.

\subsection{\texorpdfstring{4. DERIVED-UNDER-009:\allowbreak{} degree windows and relative homotopies}{4.  degree windows and relative homotopies}}\label{proofs-6209-7089-derived-under-009-degree-windows-and-relative-homotopies}

Keep I\^{}n injective and zero for n\textless a. The proof of section 2 only uses exactness of the source at degrees n>\allowbreak{}=a. It therefore proves:\allowbreak{}

\[
 H^n(C)=0\ (n\ge a)
 \quad\Longrightarrow\quad
 \text{every chain map }C\longrightarrow I
                    \text{ is null-homotopic}.          \tag{4.1}
\]

Indeed,\allowbreak{} start h\textasciicircum{}j\allowbreak{}=0 for j<\allowbreak{}=a as before. At each subsequent step n,\allowbreak{} the only source exactness used is Im(d\_\allowbreak{}C\textasciicircum{}\{n-1\})\allowbreak{}\allowbreak{}=ker(d\_\allowbreak{}C\textasciicircum{}n)\allowbreak{}. No vanishing below a is used or claimed.

Now let alpha:\allowbreak{}K \allowbreak{}-> L induce cohomology isomorphisms only in degrees n>\allowbreak{}=a. Every chain map gamma:\allowbreak{}K \allowbreak{}-> I factors uniquely through the canonical map K \allowbreak{}-> tau\_\allowbreak{}\{>\allowbreak{}=a\}K:\allowbreak{} its degree-a component kills Im(d\_\allowbreak{}K\textasciicircum{}\{a-1\})\allowbreak{},\allowbreak{} and all higher terms of the truncation are unchanged. A homotopy h\textasciicircum{}n:\allowbreak{}K\textasciicircum{}n \allowbreak{}-> I\^{}\{n-1\} has h\textasciicircum{}n\allowbreak{}=0 for n<\allowbreak{}=a,\allowbreak{} and its remaining components also factor uniquely. The homotopy equation at degree a descends because d\_\allowbreak{}K\textasciicircum{}a kills Im(d\_\allowbreak{}K\textasciicircum{}\{a-1\})\allowbreak{}. These statements give canonical,\allowbreak{} functorial isomorphisms of homotopy-class groups

\[
 \operatorname{Hom}_K(\tau_{\ge a}K,I)
       \xrightarrow{\sim}\operatorname{Hom}_K(K,I),
 \qquad
 \operatorname{Hom}_K(\tau_{\ge a}L,I)
       \xrightarrow{\sim}\operatorname{Hom}_K(L,I).        \tag{4.2}
\]

The actual map tau\_\allowbreak{}\{>\allowbreak{}=a\}(alpha)\allowbreak{} is a quasi-isomorphism. Apply section 2 to it and use (4.2)\allowbreak{}. We obtain

\[
 \alpha^*:\operatorname{Hom}_K(L,I)
             \xrightarrow{\sim}\operatorname{Hom}_K(K,I)
 \quad\text{if }H^n(\alpha)\text{ is an isomorphism for }n\ge a.
                                                               \tag{4.3}
\]

If alpha is also termwise monomorphic,\allowbreak{} the truncated map is termwise monomorphic. In its bottom degree this follows exactly from (3.1)\allowbreak{},\allowbreak{} using the injectivity of H\textasciicircum{}a(alpha)\allowbreak{}; above a it is the original monomorphism. Applying strict lifting to the truncated map and composing with L \allowbreak{}-> tau\_\allowbreak{}\{>\allowbreak{}=a\}L therefore gives,\allowbreak{} for every gamma,\allowbreak{} a strict lift beta with beta alpha\allowbreak{}=gamma.

There is a further relative assertion for two such strict lifts. Their difference factors as L \allowbreak{}-> C\allowbreak{}=L/\allowbreak{}K \allowbreak{}-> I. The long exact sequence and the isomorphisms H\textasciicircum{}n(alpha)\allowbreak{},\allowbreak{}H\textasciicircum{}\{n\allowbreak{}+1\}(alpha)\allowbreak{} give H\textasciicircum{}n(C)\allowbreak{}\allowbreak{}=0 for n>\allowbreak{}=a. Apply (4.1)\allowbreak{} to the factored map C \allowbreak{}-> I and compose its homotopy with L \allowbreak{}-> C. The resulting homotopy h satisfies

\[
 \beta_1-\beta_2=d_Ih+h d_L,\qquad h^n\alpha^n=0
                                                    \text{ for every }n. \tag{4.4}
\]

Thus both strict lifts and their uniqueness can be taken relative to the entire original subcomplex K. The source theorem only states ordinary uniqueness up to homotopy.

For precision about duality,\allowbreak{} pass from an A-complex K to the A-opposite complex DK with

\[
 (DK)^n=K^{-n},\qquad d_{DK}^n=(d_K^{-n-1})^{op}.
\]

It reverses chain maps and satisfies D(K[1]\allowbreak{})\allowbreak{}\allowbreak{}=(DK)\allowbreak{}[-1]\allowbreak{},\allowbreak{} including the minus sign in both shifted differentials. A homotopy h dualizes with component (h\textasciicircum{}\{-m\allowbreak{}+1\})\allowbreak{}\textasciicircum{}\{op\} in degree m; reversing the two compositions in its equation gives the same sum of the two required homotopy terms. Applying (4.1)\allowbreak{}–\allowbreak{}(4.4)\allowbreak{} in A-opposite proves the exact dual:\allowbreak{} if P consists of projectives with P\textasciicircum{}n\allowbreak{}=0 for n\textgreater b and alpha:\allowbreak{}L \allowbreak{}-> K induces cohomology isomorphisms for n<\allowbreak{}=b,\allowbreak{} then

\[
 \operatorname{Hom}_K(P,L)\xrightarrow{\sim}\operatorname{Hom}_K(P,K).
                                                               \tag{4.5}
\]

For termwise epimorphic alpha,\allowbreak{} any map P \allowbreak{}-> K lifts strictly. Two strict lifts beta\_\allowbreak{}1,\allowbreak{}beta\_\allowbreak{}2 have a homotopy h with alpha\textasciicircum{}\{n-1\}h\textasciicircum{}n\allowbreak{}=0 in every degree. This dual statement retains the original complexes and upper bound; it does not replace them by unspecified projective models.

\subsection{5. Projective resolutions and the cohomology bounds in lifting}\label{proofs-6209-7089-projective-resolutions-and-the-cohomology-bounds-in-lifting}

Under this exact duality,\allowbreak{} sections 1–\allowbreak{}3 prove the stated projective resolution results:\allowbreak{} bounded-above resolutions,\allowbreak{} canonical truncation tau\_\allowbreak{}\{<\allowbreak{}=b\}K with degree-b term ker(d\_\allowbreak{}K\textasciicircum{}b)\allowbreak{},\allowbreak{} termwise epimorphic resolutions for complexes whose terms vanish above b,\allowbreak{} null-homotopy of maps from bounded-above projectives to acyclic complexes,\allowbreak{} strict lifts through termwise epimorphic quasi-isomorphisms,\allowbreak{} uniqueness,\allowbreak{} and the Hom\_\allowbreak{}K to Hom\_\allowbreak{}D isomorphism.

DERIVED-RECON-041 reconciles P02-E0046 and French DERIVED-038 with the two MC-STK-ERR-0849 replacements of Comp\textasciicircum{}\allowbreak{}+ by Comp\^{}- at original 6844–\allowbreak{}6845. The plus-bounded statement would be false:\allowbreak{} let C\textasciicircum{}n\allowbreak{}=Z for n>\allowbreak{}=0,\allowbreak{} C\textasciicircum{}n\allowbreak{}=0 otherwise,\allowbreak{} and all differentials zero. The sequence 0 \allowbreak{}-> 0 \allowbreak{}-> C \allowbreak{}-> C \allowbreak{}-> 0,\allowbreak{} with the last map the identity,\allowbreak{} is a short exact sequence of bounded-below complexes. But C has nonzero cohomology in every nonnegative degree and cannot have a bounded-above projective resolution. The corrected minus-bounded statement is precisely the dual of the injective construction.

For original lemma-precise-vanishing,\allowbreak{} retain the bounded projective complex P,\allowbreak{} the term bound P\textasciicircum{}i\allowbreak{}=0 for i<n,\allowbreak{} and the cohomology bound H\textasciicircum{}i(K)\allowbreak{}\allowbreak{}=0 for i>\allowbreak{}=n. The actual canonical map tau\_\allowbreak{}\{<\allowbreak{}=n-1\}K \allowbreak{}-> K is a quasi-isomorphism. Every chain map P \allowbreak{}-> tau\_\allowbreak{}\{<\allowbreak{}=n-1\}K is zero,\allowbreak{} because in each degree one of its source or target objects is zero. The projective Hom comparison consequently gives

\[
 \operatorname{Hom}_K(P,K)=\operatorname{Hom}_D(P,K)=0.    \tag{5.1}
\]

The source also argues with the canonical truncation triangle. That triangle is in D(A)\allowbreak{}; using the same projective Hom comparison on all its vertices justifies its displayed Hom\_\allowbreak{}K groups. No termwise splitting of the canonical truncation is presumed.

For lemma-lift-map,\allowbreak{} let alpha:\allowbreak{}E \allowbreak{}-> L induce cohomology isomorphisms above n and an epimorphism in degree n,\allowbreak{} and retain the same P with terms zero below n. In the cone long exact sequence,\allowbreak{} for every i>\allowbreak{}=n the map H\textasciicircum{}i(E)\allowbreak{} \allowbreak{}-> H\textasciicircum{}i(L)\allowbreak{} is epic and H\textasciicircum{}\{i\allowbreak{}+1\}(E)\allowbreak{} \allowbreak{}-> H\textasciicircum{}\{i\allowbreak{}+1\}(L)\allowbreak{} is monic. Therefore H\textasciicircum{}i(C(alpha)\allowbreak{})\allowbreak{}\allowbreak{}=0 for i>\allowbreak{}=n. Equation (5.1)\allowbreak{} makes i beta zero in Hom\_\allowbreak{}K(P,\allowbreak{}C(alpha)\allowbreak{})\allowbreak{}. The exact Hom sequence of the original cone triangle yields gamma:\allowbreak{}P \allowbreak{}-> E with alpha gamma homotopic to beta. This proves the claimed lift with precisely the source\textquotesingle s endpoint surjectivity,\allowbreak{} rather than a stronger assumed quasi-isomorphism.

\subsection{6. Short exact sequences and DERIVED-UNDER-010}\label{proofs-6209-7089-short-exact-sequences-and-derived-under-010}

Consider 0 \allowbreak{}-> A --j-\allowbreak{}-> B --q-\allowbreak{}-> C \allowbreak{}-> 0 in bounded-below complexes. The source construction first chooses,\allowbreak{} or retains,\allowbreak{} a resolution f\_\allowbreak{}1:\allowbreak{}A \allowbreak{}-> I\_\allowbreak{}1. Its pushout is

\[
 E=\operatorname{Coker}\big((f_1,-j):A\longrightarrow I_1\oplus B\big).
                                                               \tag{6.1}
\]

It has the induced short exact sequence 0 \allowbreak{}-> I\_\allowbreak{}1 \allowbreak{}-> E \allowbreak{}-> C \allowbreak{}-> 0 and a map of short exact sequences from the original one. The outer maps are f\_\allowbreak{}1 and the identity of C,\allowbreak{} so the long exact cohomology sequences give that B \allowbreak{}-> E is a quasi-isomorphism. Since each I\_\allowbreak{}1\textasciicircum{}n is injective,\allowbreak{} the lower sequence splits in every degree. Choose the splittings and write

\[
 E^n=I_1^n\oplus C^n,\qquad
 d_E^n=\begin{pmatrix}d_{I_1}^n&\delta^n\\0&d_C^n\end{pmatrix}.
                                                               \tag{6.2}
\]

The full off-diagonal identity is d\_\allowbreak{}\{I\_\allowbreak{}1\}\textasciicircum{}\{n\allowbreak{}+1\}delta\textasciicircum{}n\allowbreak{}+delta\textasciicircum{}\{n\allowbreak{}+1\}d\_\allowbreak{}C\textasciicircum{}n\allowbreak{}=0. Thus delta:\allowbreak{}C \allowbreak{}-> I\_\allowbreak{}1[1]\allowbreak{} is a chain map with the source\textquotesingle s positive upper-right entry. This is exactly the map pi\textasciicircum{}\{n\allowbreak{}+1\}d\_\allowbreak{}E\textasciicircum{}n s\^{}n in current Homology 3606–\allowbreak{}3632; its cohomology map is the connecting map by the subsequent lemma at 3636–\allowbreak{}3648.

For the source\textquotesingle s proof,\allowbreak{} choose a termwise monomorphic resolution f\_\allowbreak{}3:\allowbreak{}C \allowbreak{}-> I\_\allowbreak{}3. Strict lifting,\allowbreak{} with target I\_\allowbreak{}1[1]\allowbreak{},\allowbreak{} gives delta':\allowbreak{}I\_\allowbreak{}3 \allowbreak{}-> I\_\allowbreak{}1[1]\allowbreak{} satisfying delta'f\_\allowbreak{}3\allowbreak{}=delta. Put

\[
 I_2^n=I_1^n\oplus I_3^n,\qquad
 d_{I_2}^n=\begin{pmatrix}d_{I_1}^n&(\delta')^n\\0&d_{I_3}^n\end{pmatrix}.
                                                               \tag{6.3}
\]

Its squared differential is zero,\allowbreak{} because the off-diagonal entry is d\_\allowbreak{}\{I\_\allowbreak{}1\}\textasciicircum{}\{n\allowbreak{}+1\}(delta')\allowbreak{}\textasciicircum{}n\allowbreak{}+(delta')\allowbreak{}\textasciicircum{}\{n\allowbreak{}+1\}d\_\allowbreak{}\{I\_\allowbreak{}3\}\textasciicircum{}n\allowbreak{}=0. The diagonal map E \allowbreak{}-> I\_\allowbreak{}2 is a chain map by delta'f\_\allowbreak{}3\allowbreak{}=delta. It is a quasi-isomorphism by the two outer quasi-isomorphisms in the short exact sequences. Composing with B \allowbreak{}-> E gives the required diagram. Every I\_\allowbreak{}2\textasciicircum{}n is injective:\allowbreak{} extension of a map to its finite biproduct is obtained by extending the two components. All complexes retain any common lower term bound,\allowbreak{} so the three source clauses,\allowbreak{} including compatibility with a prescribed I\_\allowbreak{}1,\allowbreak{} hold.

DERIVED-RECON-040 reconciles P02-E0045 and French DERIVED-037 with MC-STK-ERR-0848. The literal extra item ``add more here'' supplies no mathematical statement. Its deletion leaves exactly the three proved clauses. The strengthening below is not inserted in place of that placeholder.

\textbf{DERIVED-UNDER-010.} Both outer resolutions can in fact be prescribed:\allowbreak{} take any given f\_\allowbreak{}1:\allowbreak{}A \allowbreak{}-> I\_\allowbreak{}1 and f\_\allowbreak{}3:\allowbreak{}C \allowbreak{}-> I\_\allowbreak{}3 satisfying the source definition,\allowbreak{} with no monomorphism requirement on f\_\allowbreak{}3. Once these resolutions exist,\allowbreak{} no additional enough-injectives hypothesis is needed. Keep (6.1)\allowbreak{}–\allowbreak{}(6.2)\allowbreak{}. Ordinary lifting up to homotopy gives delta':\allowbreak{}I\_\allowbreak{}3 \allowbreak{}-> I\_\allowbreak{}1[1]\allowbreak{} and maps h\textasciicircum{}n:\allowbreak{}C\textasciicircum{}n \allowbreak{}-> I\_\allowbreak{}1\textasciicircum{}n with the exact equation

\[
 (\delta')^n f_3^n-\delta^n
       =-d_{I_1}^n h^n+h^{n+1}d_C^n.                    \tag{6.4}
\]

The minus sign is the one in d\_\allowbreak{}\{I\_\allowbreak{}1[1]\allowbreak{}\}. Define I\_\allowbreak{}2 by (6.3)\allowbreak{},\allowbreak{} but now use the map

\[
 e^n:E^n\longrightarrow I_2^n,\qquad
 e^n=\begin{pmatrix}1&h^n\\0&f_3^n\end{pmatrix}.          \tag{6.5}
\]

The two possible upper-right entries in its chain-map equation are

\[
 (d_{I_2}^n e^n)_{12}=d_{I_1}^n h^n+(\delta')^n f_3^n,\qquad
 (e^{n+1}d_E^n)_{12}=\delta^n+h^{n+1}d_C^n.
\]

They are equal by (6.4)\allowbreak{}; the other three entries agree by identity and the chain equation for f\_\allowbreak{}3. Thus e is an actual chain map. Its restriction on I\_\allowbreak{}1 is the identity,\allowbreak{} and its map on the quotient is exactly the prescribed f\_\allowbreak{}3. The diagram

\[
\begin{array}{ccccccccc}
0&\longrightarrow&A&\longrightarrow&B&\longrightarrow&C&\longrightarrow&0\\
&&\downarrow f_1&&\downarrow&&\downarrow f_3&&\\
0&\longrightarrow&I_1&\longrightarrow&I_2&\longrightarrow&I_3&\longrightarrow&0
\end{array}                                                    \tag{6.6}
\]

therefore commutes strictly,\allowbreak{} its rows are exact,\allowbreak{} and the middle vertical map is a quasi-isomorphism. The lower sequence is termwise split. No term of I\_\allowbreak{}1 or I\_\allowbreak{}3 has been changed. A common lower bound on the original terms and the prescribed resolutions is retained.

For completeness the projective statement has an equally explicit construction. Prescribe f\_\allowbreak{}1:\allowbreak{}P\_\allowbreak{}1 \allowbreak{}-> A and f\_\allowbreak{}3:\allowbreak{}P\_\allowbreak{}3 \allowbreak{}-> C,\allowbreak{} and form E\allowbreak{}=B times\_\allowbreak{}C P\_\allowbreak{}3. The map E \allowbreak{}-> B is a quasi-isomorphism,\allowbreak{} and 0 \allowbreak{}-> A \allowbreak{}-> E \allowbreak{}-> P\_\allowbreak{}3 \allowbreak{}-> 0 is termwise split. Write d\_\allowbreak{}E\allowbreak{}=[[d\_\allowbreak{}A,\allowbreak{}delta]\allowbreak{},\allowbreak{}[0,\allowbreak{}d\_\allowbreak{}\{P\_\allowbreak{}3\}]\allowbreak{}]\allowbreak{},\allowbreak{} with delta:\allowbreak{}P\_\allowbreak{}3 \allowbreak{}-> A[1]\allowbreak{}. Projective lifting through f\_\allowbreak{}1[1]\allowbreak{} supplies delta':\allowbreak{}P\_\allowbreak{}3 \allowbreak{}-> P\_\allowbreak{}1[1]\allowbreak{} and h\textasciicircum{}n:\allowbreak{}P\_\allowbreak{}3\textasciicircum{}n \allowbreak{}-> A\^{}n such that

\[
 \delta^n-f_1^{n+1}(\delta')^n
              =-d_A^n h^n+h^{n+1}d_{P_3}^n.
\]

Set d\_\allowbreak{}\{P\_\allowbreak{}2\}\allowbreak{}=[[d\_\allowbreak{}\{P\_\allowbreak{}1\},\allowbreak{}delta']\allowbreak{},\allowbreak{}[0,\allowbreak{}d\_\allowbreak{}\{P\_\allowbreak{}3\}]\allowbreak{}]\allowbreak{} on P\_\allowbreak{}1 direct-sum P\_\allowbreak{}3 and use the map P\_\allowbreak{}2 \allowbreak{}-> E with matrix [[f\_\allowbreak{}1,\allowbreak{}h]\allowbreak{},\allowbreak{}[0,\allowbreak{}1]\allowbreak{}]\allowbreak{}. Its upper-right chain equation is d\_\allowbreak{}A h\allowbreak{}+delta\allowbreak{}=f\_\allowbreak{}1[1]\allowbreak{}delta'\allowbreak{}+h d\_\allowbreak{}\{P\_\allowbreak{}3\},\allowbreak{} exactly the displayed equation; the remaining entries are the original chain equations. Composing with E \allowbreak{}-> B gives a short exact sequence of projective resolutions with both outer resolutions fixed. Finite biproducts of projectives are projective,\allowbreak{} by lifting the two source components. Any common upper bound is retained. This also proves the projective receiving version without hiding the dual signs.

\subsection{7. Computing right derived functors and higher connecting maps}\label{proofs-6209-7089-computing-right-derived-functors-and-higher-connecting-maps}

For bounded-below injective I,\allowbreak{} let S\_\allowbreak{}I be the source\textquotesingle s index category I/\allowbreak{}Qis\textasciicircum{}\allowbreak{}+(A)\allowbreak{},\allowbreak{} formed in the homotopy category. Every object alpha:\allowbreak{}I \allowbreak{}-> K has a morphism to the identity object:\allowbreak{} a map beta:\allowbreak{}K \allowbreak{}-> I with beta alpha\allowbreak{}=1 in K(A)\allowbreak{},\allowbreak{} obtained in section 2. It is a quasi-isomorphism,\allowbreak{} since alpha is. This beta is unique in K(A)\allowbreak{},\allowbreak{} again by the isomorphism alpha\^{}* of section 2. Thus the identity index is terminal,\allowbreak{} and its singleton subcategory is cofinal. For any exact F:\allowbreak{}K\textasciicircum{}\allowbreak{}+(A)\allowbreak{} \allowbreak{}-> D the diagram of target values F(K)\allowbreak{} is consequently essentially constant with value F(I)\allowbreak{}. The canonical comparison at the identity is the identity of F(I)\allowbreak{}; this proves that I computes RF. The argument does not need enough injectives in A.

Taking I to be an injective object in degree zero,\allowbreak{} and taking the exact complex-level functor induced by an additive A \allowbreak{}-> B,\allowbreak{} proves that this object is right acyclic. The passage between bounded complex-level and abelian-category derived functors is the exact bounded comparison proved in section 2 of PROOFS\_\allowbreak{}5091\_\allowbreak{}5900.md. DERIVED-RECON-042 reconciles P02-E0047 and French DERIVED-039 with MC-STK-ERR-0850\textquotesingle s singular ``direct consequence'' at original 6962.

With enough injectives,\allowbreak{} every bounded-below complex has an injective resolution. Its computing object just proved gives existence of RF on every object of D\textasciicircum{}\allowbreak{}+(A)\allowbreak{},\allowbreak{} for any exact F from K\textasciicircum{}\allowbreak{}+(A)\allowbreak{}. For an additive A \allowbreak{}-> B this also follows from enough right acyclic objects. That latter route uses the repaired proposition-enough-acyclics:\allowbreak{} an acyclic complex is a zero object in the derived category,\allowbreak{} so RF is individually defined there,\allowbreak{} and the preceding Leray theorem proves its F-image acyclic. The counterexample to the old image induction and the noncircular replacement are fully proved in PROOFS\_\allowbreak{}5091\_\allowbreak{}5900.md,\allowbreak{} section 6. The present existence argument therefore does not reinstate the erroneous image-closure assertion.

The exactness of the everywhere defined derived functor,\allowbreak{} including its original shift comparison,\allowbreak{} was proved by the two-out-of-three construction in PROOFS\_\allowbreak{}4166\_\allowbreak{}5089.md. Compose it with the exact K\textasciicircum{}\allowbreak{}+(A)\allowbreak{} \allowbreak{}-> D\textasciicircum{}\allowbreak{}+(A)\allowbreak{} to obtain the stated exact functor on K\textasciicircum{}\allowbreak{}+(A)\allowbreak{}. For a short exact sequence of complexes,\allowbreak{} the canonical derived connecting triangle gives,\allowbreak{} after RF,\allowbreak{} a distinguished triangle; this is the asserted delta-functor on Comp\textasciicircum{}\allowbreak{}+(A)\allowbreak{}. Its restriction along A \allowbreak{}-> Comp\textasciicircum{}\allowbreak{}+(A)\allowbreak{},\allowbreak{} A mapped to A[0]\allowbreak{},\allowbreak{} is again a delta-functor. Alternatively,\allowbreak{} (6.6)\allowbreak{} computes that triangle by a termwise split sequence,\allowbreak{} whose F-image is still termwise split for any additive F.

For the source\textquotesingle s left exact F,\allowbreak{} the long exact cohomology sequence of this triangle is precisely the displayed sequence at 7057–\allowbreak{}7061. Negative R\^{}iF vanish by the lower-bound argument,\allowbreak{} and the canonical F \allowbreak{}-> R\^{}0F is an isomorphism by the kernel calculation in PROOFS\_\allowbreak{}5091\_\allowbreak{}5900.md,\allowbreak{} section 4. Positive R\^{}iF vanish on injectives because they are computed in degree zero. The actual connecting map R\textasciicircum{}iF(C)\allowbreak{} \allowbreak{}-> R\textasciicircum{}\{i\allowbreak{}+1\}F(A)\allowbreak{} is H\^{}i of RF of the canonical C \allowbreak{}-> A[1]\allowbreak{},\allowbreak{} followed by the inverse of the exact-functor shift comparison. This retains the original connecting-map sign.

Every A embeds into an injective,\allowbreak{} so all positive R\^{}iF are effaceable. The complete effacement-to-universality proof is in ../\allowbreak{}HOMOLOGY\_\allowbreak{}INTAKE\_\allowbreak{}20260930/\allowbreak{}PROOFS\_\allowbreak{}1901\_\allowbreak{}2765.md,\allowbreak{} lines 310–\allowbreak{}374,\allowbreak{} with its application to the present canonical connecting maps in PROOFS\_\allowbreak{}5091\_\allowbreak{}5900.md,\allowbreak{} section 4. Specifically R\textasciicircum{}1F(A)\allowbreak{} is the cokernel of F(I)\allowbreak{} \allowbreak{}-> F(I/\allowbreak{}A)\allowbreak{},\allowbreak{} and higher R\textasciicircum{}\{n\allowbreak{}+1\}F(A)\allowbreak{} identify with R\textasciicircum{}nF(I/\allowbreak{}A)\allowbreak{}; these maps construct and uniquely determine the extension of a degree-zero natural transformation. Comparing embeddings via their common injective direct sum gives independence,\allowbreak{} and naturality of the short exact sequences gives the delta-functor equations. All effacement hypotheses are thus realized by the actual injective embeddings,\allowbreak{} not assumed as unproved conditions.

\subsection{8. Receiving maps and propagation}\label{proofs-6209-7089-receiving-maps-and-propagation}

The correction NEW-024 preserves the original lifting and uniqueness theorems,\allowbreak{} (3.6)\allowbreak{},\allowbreak{} and their projective duals. Their cited uses therefore receive corrected proofs without changed theorem hypotheses. The following bounded checks identify the actual receiving maps.

In current Cohomology 2440–\allowbreak{}2464,\allowbreak{} G[0]\allowbreak{} \allowbreak{}-> J is an object injective resolution and hence a termwise monomorphism. The map G[0]\allowbreak{} \allowbreak{}-> f\_\allowbreak{}*I \allowbreak{}-> J\textquotesingle{} has the bounded-below injective target J\textquotesingle. Strict lifting gives the displayed beta:\allowbreak{}J \allowbreak{}-> J\textquotesingle. Formula (4.4)\allowbreak{} shows that two such strict choices are homotopic by a homotopy zero on G[0]\allowbreak{}. Applying the additive global-section functor preserves its two-term homotopy equation. The stated cohomology map is unchanged. In current Descent 3933–\allowbreak{}3956 the same construction,\allowbreak{} with source F[0]\allowbreak{} \allowbreak{}-> J|\_\allowbreak{}S and injective target I,\allowbreak{} gives the map J|\_\allowbreak{}S \allowbreak{}-> I used there. This check uses the exactness established in that receiving passage; it does not claim a new audit of its preceding vanishing theorems.

For the flat base-change constructions in current Cohomology 3368–\allowbreak{}3413 and Cohomology on Sites 2439–\allowbreak{}2484,\allowbreak{} the target (g')\allowbreak{}\_\allowbreak{}*J is a bounded-below complex of injectives. Indeed,\allowbreak{} for a monomorphism M \allowbreak{}-> N,\allowbreak{} exactness of (g')\allowbreak{}\textasciicircum{}* and injectivity of J\^{}n make Hom((g')\allowbreak{}\textasciicircum{}*N,\allowbreak{}J\textasciicircum{}n)\allowbreak{} \allowbreak{}-> Hom((g')\allowbreak{}\textasciicircum{}*M,\allowbreak{}J\textasciicircum{}n)\allowbreak{} epic; adjunction proves that (g')\allowbreak{}\_\allowbreak{}*J\textasciicircum{}n is injective. Sections 2–\allowbreak{}3 therefore construct the actual beta:\allowbreak{}I \allowbreak{}-> (g')\allowbreak{}\_\allowbreak{}*J and prove uniqueness up to homotopy. The map f\_\allowbreak{}*beta is carried by the original identification f\_\allowbreak{}*(g')\allowbreak{}\_\allowbreak{}* \allowbreak{}= g\_\allowbreak{}*(f')\allowbreak{}\_\allowbreak{}* and the adjunction to g\textasciicircum{}*f\_\allowbreak{}*I \allowbreak{}-> (f')\allowbreak{}\_\allowbreak{}*J. Additivity of f\_\allowbreak{}* preserves the homotopy equation,\allowbreak{} and the adjunction bijection in each degree preserves compositions with the differentials and addition. Thus the homotopy descends to the displayed base-change arrow. Comparing different injective resolutions by (3.6)\allowbreak{} and the same uniqueness argument gives the same derived arrow under their canonical resolution identifications. No new flatness or source comparison hypothesis is introduced.

UNDER-009 applies to these comparisons with the actual lower term bound of (g')\allowbreak{}\_\allowbreak{}*J:\allowbreak{} cohomological agreement in those degrees already suffices for the Hom comparison. This is recorded as a possible choice of comparison map justified by (4.3)\allowbreak{},\allowbreak{} while the source\textquotesingle s full quasi-isomorphisms remain in its proof.

For the Cech boundary comparison in current Cohomology 3150–\allowbreak{}3204,\allowbreak{} UNDER-010 allows any already chosen outer injective resolutions I\_\allowbreak{}1 and I\_\allowbreak{}3 to be retained in the displayed short exact sequence. The actual I\_\allowbreak{}2 of (6.3)\allowbreak{} and map (6.5)\allowbreak{} supply its middle term. Applying sections on every intersection of the original covering preserves the split sequence. Taking the stated products over intersection indices preserves this splitting as well,\allowbreak{} with the entire indexing family retained. The total-complex maps consequently form the same exact diagram. The connecting map is induced by the actual off-diagonal delta\textquotesingle{} from (6.3)\allowbreak{},\allowbreak{} with the sign verified in Homology 3606–\allowbreak{}3648. Naturality of the exact diagram under refinement proves the compatibility of the two boundary constructions as in the receiving proof. The cover and both outer resolutions are unchanged.

In current More Algebra 16590–\allowbreak{}16626,\allowbreak{} P is the bounded finite-free complex with P\textasciicircum{}i\allowbreak{}=0 for i<m\allowbreak{}+1,\allowbreak{} and E \allowbreak{}-> L is an isomorphism on cohomology above m and surjective in degree m. Represent the derived composite P \allowbreak{}-> K \allowbreak{}-> L by a chain map,\allowbreak{} using the projective Hom comparison. Section 5 applies with n\allowbreak{}=m\allowbreak{}+1 and gives the required gamma:\allowbreak{}P \allowbreak{}-> E. Its cone is the original complex of finite free modules,\allowbreak{} with terms P\textasciicircum{}\{i\allowbreak{}+1\} direct-sum E\^{}i and differential (-d\_\allowbreak{}P,\allowbreak{}0;gamma,\allowbreak{}d\_\allowbreak{}E)\allowbreak{}; it remains bounded.

For the cohomology bound on C(gamma)\allowbreak{} \allowbreak{}-> M,\allowbreak{} retain the two exact triangle sequences. In degree i>m\allowbreak{}+1,\allowbreak{} the adjacent maps for P \allowbreak{}-> K and E \allowbreak{}-> L are isomorphisms,\allowbreak{} so the cone map is an isomorphism. At i\allowbreak{}=m\allowbreak{}+1,\allowbreak{} P \allowbreak{}-> K is still surjective at the left endpoint and is an isomorphism in degree m\allowbreak{}+2,\allowbreak{} while E \allowbreak{}-> L is an isomorphism in both relevant degrees. A class in the kernel of the cone map has zero boundary in H\textasciicircum{}\{m\allowbreak{}+2\}(P)\allowbreak{},\allowbreak{} hence comes from H\textasciicircum{}\{m\allowbreak{}+1\}(E)\allowbreak{}; its image in H\textasciicircum{}\{m\allowbreak{}+1\}(L)\allowbreak{} comes from H\textasciicircum{}\{m\allowbreak{}+1\}(K)\allowbreak{},\allowbreak{} which lifts to H\textasciicircum{}\{m\allowbreak{}+1\}(P)\allowbreak{}. Subtracting that image makes the class zero,\allowbreak{} proving injectivity. Surjectivity follows by the same exact sequence,\allowbreak{} with the two next endpoint isomorphisms. At i\allowbreak{}=m,\allowbreak{} lift the boundary of a class of H\textasciicircum{}m(M)\allowbreak{} to H\textasciicircum{}\{m\allowbreak{}+1\}(P)\allowbreak{},\allowbreak{} using surjectivity there. Its image in H\textasciicircum{}\{m\allowbreak{}+1\}(E)\allowbreak{} is zero,\allowbreak{} since that group maps isomorphically to H\textasciicircum{}\{m\allowbreak{}+1\}(L)\allowbreak{}. Lift it to H\textasciicircum{}m(C(gamma)\allowbreak{})\allowbreak{}; the difference of its image from the original class comes from H\textasciicircum{}m(L)\allowbreak{},\allowbreak{} and can be removed using the epimorphism H\textasciicircum{}m(E)\allowbreak{} \allowbreak{}-> H\textasciicircum{}m(L)\allowbreak{}. Thus the cone map is surjective in degree m. This proves exactly the receiving pseudo-coherence estimate,\allowbreak{} including both endpoint distinctions.

These are dependency checks on the listed passages,\allowbreak{} not a review of the whole receiving chapters. No source operation in those chapters is needed for these results.
